\ProvidesFile{Tuple.tex}[Tuple]


\subsection{\texttt{Tuple}}
\label{section::Tuple}

\subsubsection{Зачем?}

Учитывая, что STL вполне себе содержит тип \verb"std::tuple", надо бы объяснить а зачем мне нужно писать свой.
Главная для меня проблема заключается в том, что  \verb"std::tuple" невозможно использовать в \verb"constexpr" контексте, а именно, я не могу сделать следующее:
\begin{cppcode}[\nonumbers]
template<class T, T x>
struct Value {};

using Data = Value<std::tuple<int, double, float>, std::make_tuple<int, double, float>(1)>;
\end{cppcode}
Причина в том, что в \verb"std::tuple" есть приватные данные, а элементы таких классов нельзя использовать в качестве параметров шаблона.
Другая проблема в том, что \verb"std::tuple" имеет обращение к элементу за \verb"O(n)", что так себе.
А потому цель этого раздела в том, чтобы имплементировать \verb"constexpr" версию \verb"Tuple", которую я буду использовать для имплементации State Machine позже.%
\footnote{С появлением оператора \texttt{operator[]} для аргумент pack-а в C++26 многое из того, что здесь сказано можно сделать чуть проще.}

\subsubsection{Имплементация}

Основная идея заключается в том, чтобы не определять \verb"Tuple" индуктивно, вместо этого под каждый тип заведем отдельный класс \verb"Node<ind, T>" и унаследуемся от этих нодов, где \verb"ind" -- целочисленный индекс номера элемента в \verb"Tuple".
Наличие индекса различает элементы даже при одинаковом \verb"T", что не накладывает дополнительных ограничений на типы.
Кроме того, это позволит получать доступ к элементу с помощью \verb"static_cast" трюка за \verb"O(1)".
Начнем с описания \verb"Node":
\begin{cppcode}
namespace detail {
template<int ind, class T>
struct Node {
  using value_type = int;
  static constexpr value_type value = ind;
  using type = T;
  type data;
};
} // namespace detail
\end{cppcode}
В саму структуру зашьем индекс и тип, а так же хранимый элемент \verb"data".
Теперь определим приватную имплементацию для \verb"Tuple" следующим образом:
\begin{cppcode}
template<class List, class... Ts>
struct TupleImpl {
  static_assert(Logic::False<List>, "Internal error of Tuple!");
};

template<template<class U, U...> class List, int... inds, class... Ts>
struct TupleImpl<List<int, inds...>, Ts...> : Node<inds, Ts>... {
  /*ctors*/
};
\end{cppcode}
Имплементация зависит от списка \verb"List", который содержит последовательность целых чисел вида \verb"EL::List<int, 0, 1,..., n-1>" и набора типов \verb"Ts...".
При этом предполагается, что \verb"n" совпадает с количеством элементов в \verb"Ts...".
Для нас основным будет специализация, где список типов раскрывается, а общее определение используется для сообщения о внутренней ошибке имплементации (первая общая версия не должна никогда вызываться).

Я пока пропустил код для двух конструкторов, которые мне понадобятся, чтобы они не мешали.
Чуть позже я их добавлю.
А теперь давайте я возьму пример со списком из трех типов:
\begin{cppcode}[\nonumbers]
struct Tuple<List<int, 0, 1, 2>, int, bool, char> : Node<0, int>, Node<1, bool>, Node<2, char> {
  /*ctors*/
};
\end{cppcode}
Так как в коде используются структуры с публичными полями, то достучаться до конкретного \verb"Node" можно с помощью \verb"static_cast" вот как:
\begin{cppcode}[\nonumbers]
using CTuple = Tuple<List<int, 0, 1, 2>, int, bool, char>;

CTuple t{1, true, 'c'};
int = static_cast<Node<0, int>&>(t).data;
bool = static_cast<Node<1, bool>&>(t).data;
char = static_cast<Node<2, int>&>(t).data;
\end{cppcode}
Я чуть позже упрощу этот код с помощью функции \verb"Get", а пока пусть это будет лишь демонстрация устройства самого типа.
Теперь добавим конструкторы:
\begin{cppcode}
template<template<class U, U...> class List, int... inds, class... Ts>
struct TupleImpl<List<int, inds...>, Ts...> : Node<inds, Ts>... {
  constexpr TupleImpl() = default;
  
  template<class... Args, class = std::enable_if_t<(sizeof...(Args) == sizeof...(Ts))>>
  constexpr TupleImpl(Args&&... args)
      : Node<inds, Ts>{std::forward<Ts>(args)}... {
  }
};
\end{cppcode}
Дефолтный конструктор просто по умолчанию и обязательно \verb"constexpr".
А во втором конструкторе обратите внимание на то, что мы используем шаблонный конструктор
\begin{cppcode}
template<class... Args, class = std::enable_if_t<(sizeof...(Args) == sizeof...(Ts))>>
constexpr TupleImpl(Args&&... args)
\end{cppcode}
Так приходится делать, если мы хотим поддержать perfect forwarding (см.~\ref{section::InitClass} пункт~\ref{section::InitClass::PerfForward}).
Если бы мы написали 
\begin{cppcode}
template<class List, class... Ts>
struct TupleImpl {
  static_assert(Logic::False<List>, "Internal error of Tuple!");
};

template<template<class U, U...> class List, int... inds, class... Ts>
struct TupleImpl<List<int, inds...>, Ts...> : Node<inds, Ts>... {
  constexpr TupleImpl(Ts&&... args)
      : Node<inds, Ts>{std::forward<Ts>(args)}... {
  }
};
\end{cppcode}
То аргументы уже были бы rvalue ссылками, так как теперь внутри \verb"TupleImpl" типы \verb"Ts..." известные, а не шаблонные параметры.
А мы хотим передавать не только rvalue ссылки, например такой конструктор нельзя забиндить на переменные, потому что они lvalue.
По хорошему надо все возможные сочетания \verb"&&" и \verb"&" по всем аргументам.
Но это как раз делается с помощью perfect forwarding.
А для этого типы аргументов должны быть шаблонными параметрами.
И мы делаем только одну простую проверку, что количество типов такое же, что и количество аргументов.
А дальше уже попытка присвоить аргумент из \verb"args" к полю \verb"data" внутри \verb"Node", которое происходит тут
\begin{cppcode}[\nonumbers]
Node<inds, Ts>{std::forward<Ts>(args)}...
\end{cppcode}
вызовет ошибку, при невозможности присвоить.
Теперь можно определить сам шаблон \verb"Tuple" так:
\begin{cppcode}
template<class... Ts>
struct Tuple : detail::TupleImpl<EL::MakeList<int, 0, sizeof...(Ts)>, Ts...> {
  using Base = detail::TupleImpl<EL::MakeList<int, 0, sizeof...(Ts)>, Ts...>;
  using Base::Base;
  
  static constexpr int size = sizeof...(Ts);
};
\end{cppcode}
Здесь мы просто наследуемся публично от \verb"TupleImpl" и передаем ему сгенерированную последовательность из целых индексов для всех элементов в \verb"Ts..." (см.~\ref{section::EListImpl}) и список типов \verb"Ts...".
Так же мы наследуем все конструкторы в $4$ строчке.
И вводим удобную константу \verb"size" для получения размера \verb"Tuple" .

\subsubsection{Проверки и доступ}

\paragraph{\texttt{isTuple}}

Если мы хотим работать с новым типом, то для начала надо научиться проверять он ли нам дан:
\begin{cppcode}
namespace detail {
template<class T>
struct isTupleImpl : Bool_t<false> {};

template<class... Ts>
struct isTupleImpl<Tuple<Ts...>> : Bool_t<true> {};
} // namespace detail

template<class T>
struct isTuple_t : detail::isTupleImpl<std::remove_cv_t<T>> {};

template<class T>
constexpr bool isTuple = isTuple_t<T>::value;
\end{cppcode}

\paragraph{\texttt{Get}}

Для получения доступа к элементу по индексу хорошо бы сделать красивый метод обертку.
В начале определим функцию без каких-либо проверок:
\begin{cppcode}
namespace detail {
template<int ind, class T>
constexpr T& Get_Impl(detail::Node<ind, T>& t) noexcept {
  return t.data;
}

template<int ind, class T>
constexpr const T& Get_Impl(const detail::Node<ind, T>& t) noexcept {
  return t.data;
}

template<int ind, class T>
constexpr T&& Get_Impl(detail::Node<ind, T>&& t) noexcept {
  return std::move(t).data;
}

template<int ind, class T>
constexpr const T&& Get_Impl(const detail::Node<ind, T>&& t) noexcept {
  return std::move(t).data;
}
} // namespace detail
\end{cppcode}
Возникает вопрос, какая у них должна быть сигнатура.
И правильный ответ надо смотреть у метода \verb"std::get", который определяет ровно эти $4$ сигнатуры.
Теперь можно определить более безопасные версии с проверками:
\begin{cppcode}
template<int ind, class... Ts>
constexpr auto& Get(Tuple<Ts...>& t) noexcept {
  static_assert(0 <= ind, "The index of Get MUST be non-negative!");
  static_assert(ind < Tuple<Ts...>::size, "The index of Get MUST be in range of the Tuple!");
  return detail::Get_Impl<ind>(t);
}

template<int ind, class... Ts>
constexpr const auto& Get(const Tuple<Ts...>& t) noexcept {
  static_assert(0 <= ind, "The index of Get MUST be non-negative!");
  static_assert(ind < Tuple<Ts...>::size, "The index of Get MUST be in range of the Tuple!");
  return detail::Get_Impl<ind>(t);
}

template<int ind, class... Ts>
constexpr auto&& Get(Tuple<Ts...>&& t) noexcept {
  static_assert(0 <= ind, "The index of Get MUST be non-negative!");
  static_assert(ind < Tuple<Ts...>::size, "The index of Get MUST be in range of the Tuple!");
  return detail::Get_Impl<ind>(std::move(t));
}

template<int ind, class... Ts>
constexpr const auto&& Get(const Tuple<Ts...>&& t) noexcept {
  static_assert(0 <= ind, "The index of Get MUST be non-negative!");
  static_assert(ind < Tuple<Ts...>::size, "The index of Get MUST be in range of the Tuple!");
  return detail::Get_Impl<ind>(std::move(t));
}
\end{cppcode}
Аналогично можно сделать доступ по имени типа.
Я приведу пример только для не константной rvalue ссылки, чтобы не писать груду одинакового кода.
\begin{cppcode}
namespace detail {
template<class T, int ind>
constexpr T&& Get_Impl(Node<ind, T>&& t) noexcept {
  return std::move(t).data;
}
} // namespace detail

template<class T, class... Ts>
constexpr T&& Get(Tuple<Ts...>&& t) noexcept {
  static_assert(TL::Contains<TL::List<Ts...>, T>, "The type T MUST be in the tuple!");
  return detail::Get_Impl<T>(std::move(t));
}
\end{cppcode}
Теперь можно получать доступ к элементам и по индексу и по имени типа:
\begin{cppcode}
using Data = Tuple<int, bool, char>;
Data t{2, true, 'c'};

int  x = Get<0>(t);
bool y = Get<1>(t);
char z = Get<2>(t);

int  u = Get<int>(t);
bool v = Get<bool>(t);
char w = Get<char>(t);
\end{cppcode}

\paragraph{\texttt{TypeAt}}

Теперь сделаем доступ к типу элемента по индексу.
Как обычно снаружи сделаем удобный интерфейс для пользования:
\begin{cppcode}
template<class T, int ind>
struct TypeAt_t : detail::TypeAtImpl<T, ind> {};

template<class T, int ind>
using TypeAt = typename TypeAt_t<T, ind>::type;
\end{cppcode}
И всю содержательную работу спрячем в \verb"TypeAtImpl".
Я как обычно в имплементации ввожу дополнительный флаг и использую общий шаблон с ложным флагом для сообщении об ошибках и специализацию для \verb"true" флага уже как рабочую имплементацию.
\begin{cppcode}
template<class T, int ind, bool isCorrect = TypeAtCheck<T, ind>>
struct TypeAtImpl {
  static_assert(isTuple<T>, "The first argument of TypeAt MUST be a Tuple!");
  static_assert(ind >= 0, "The second argument of TypeAt MUST be non-negative!");
  static_assert(Imply_u<isTuple_t<T>, isLess_w<Int_t<ind>, TL::Size_t<T>>>,
               "The second argument of TypeAt MUST be less than the amount of types in Tuple!");
  using type = void;
};
\end{cppcode}
Функции проверки выглядят так
\begin{cppcode}
template<class T, int ind>
struct TypeAtCheck_t {
  using IsTuple = isTuple_t<T>;
  using IsNonNegative = Bool_t<ind >= 0>;
  using IsBounded = isLess_w<Int_t<ind>, TL::Size_t<T>>;
  using value_type = bool;
  static constexpr value_type value = And_u<IsTuple, IsNonNegative, IsBounded>;
};

template<class T, int ind>
constexpr bool TypeAtCheck = TypeAtCheck_t<T, ind>::value;
\end{cppcode}
Здесь я пользуюсь ленивым вычислением \verb"And" шаблонов.
Если первое условие не верное, то второе и третье уже не проверяются.
А потому если мы проверили, что шаблон является \verb"Tuple", то вычисление его размера с помощью \verb"TL::Size" не упадет (либо можно было воспользоваться полем \verb"size", но тогда его бы пришлось заворачивать во вспомогательный шаблон в силу особенностей компиляции шаблонного кода).
И теперь код для собственно имплементации.
В начале будет вспомогательная функция и потом уже сама имплементация.
Я опять воспользуюсь тем, что функции могут восстанавливать пропущенные шаблонные аргументы:
\begin{cppcode}
template<int ind, class T>
consteval T Type_At_Impl(detail::Node<ind, T> t) noexcept {
  return t.data;
}

template<class T, int ind>
struct TypeAtImpl<T, ind, true> {
  using type = decltype(Type_At_Impl<ind>(std::declval<T>()));
};
\end{cppcode}
То есть вспомогательная функция мне возвращает объект типа \verb"T, который лежит в \verb"Tuple", а потом я уже с помощью \verb"decltype" перевожу элемент в тип.
Тут воспользоваться \verb"Get" нельзя, так как он возвращает ссылки, а нам нужен в точности тот тип, который лежит внутри.

\subsubsection{Forward и ForwardAsTuple}

Если уж мы начали имплементировать свой \verb"Tuple" и переходить на него с \verb"std::tuple", то мы никуда не денемся без имплементации своей версии \verb"ForwardAsTuple", а значит нужно и обычную версию \verb"Forward" написать.
Обычная \verb"Forward" пишется так:
\begin{cppcode}
template<class T>
constexpr T&& Forward(std::remove_reference_t<T>& t) noexcept {
  return static_cast<T&&>(t);
}

template<class T>
constexpr T&& Forward(std::remove_reference_t<T>&& t) noexcept {
  return static_cast<T&&>(t);
}
\end{cppcode}
А \verb"ForwardAsTuple" тогда имплементируется следующим образом:
\begin{cppcode}
template<class... Args>
auto ForwardAsTuple(Args&&... args) {
  using T = Tuple<Args&&...>;
  return T(Forward<Args>(args)...);
}
\end{cppcode}

\subsubsection{Apply}

Работая над списками типов и элементов в compile-time, мы научились применять различные операции к элементам списков с помощью \verb"Map" конструкции.
Здесь же у нас \verb"Tuple" хранит информацию в виде нестатических членов, а потому конструкция \verb"Map" не сработает.
А потому надо определить ее аналог.
Я не знаю зачем я ее назвал \verb"Apply", ну пусть так и будет.
Теперь как выглядит сама задача.
Пусть у нас есть некоторый шаблонный класс вида:
\begin{cppcode}
template<auto x>
struct F {
  using value_type = /*some type*/;
  static constexpr value_type value = /*some formula*/;
};
\end{cppcode}
Тогда \verb"F<x>::value" представляет из себя результат применения \verb"F" к элементу \verb"x".
Пусть теперь нам дан \verb"Tuple" из каких-то элементов и наша задача построить новый \verb"Tuple" применив \verb"F" к каждому элементу исходного \verb"Tuple".
Имплементацию делаем по стандартной схеме.
Публичная часть выглядит так:
\begin{cppcode}
template<template<auto> class F, auto t>
using Apply_t = typename detail::ApplyImpl<F, std::remove_cv_t<decltype(t)>, t>::type;

template<template<auto> class F, auto t>
constexpr auto Apply = detail::ApplyImpl<F, std::remove_cv_t<decltype(t)>, t>::value;
\end{cppcode}
Здесь должно броситься в глаза, что я передаю вместе с элементом \verb"t" сразу и его тип и убираю с него \verb"const/volatile" спецификаторы.
Еще я определил нестандартным для себя образом \verb"Apply_t", и потому \verb"constexpr" переменная \verb"Apply" не получается из \verb"Apply_t" распаковкой.
Просто мне нужна версия, которая действует и на типах и на элементах.
Я когда-нибудь причешу этот код.
Как и раньше я использую общую версию шаблона для сообщения об ошибках: 
\begin{cppcode}
template<template<auto> class F, class T, auto t>
struct ApplyImpl {
  static_assert(TL::isList<T>, "The second element of Apply MUST be of Tuple-like type!");
};
\end{cppcode}
При этом я не использую никакую проверку в дополнительном флаге.
Я просто распакую \verb"Tuple".
Причина в том, что у меня нет нормальной проверки, что мне дали тип похожий на \verb"Tuple".
Ее можно написать, но я решил сократить.
\begin{cppcode}
template<template<auto> class F, class... Ts, template<class...> class Tuple, Tuple<Ts...> t>
struct ApplyImpl<F, Tuple<Ts...>, t> {
  using Indices = EL::MakeList<int, 0, TL::Size<Tuple<Ts...>>>;
  static constexpr auto value = ApplyByIndices<F, Indices, Tuple, t>;
  using value_type = std::remove_cv_t<decltype(value)>;
  using type = value_type;
};
\end{cppcode}
Внутри имплементации я создаю последовательность индексов для всех элементов \verb"Indices" и передаю вычисление другому шаблону \verb"ApplyByIndices", который должен будет ее распаковать.
Его имплементация выглядит так:
\begin{cppcode}
template<template<auto> class F, class Ind, template<class...> class Tuple, auto t>
struct ApplyByIndicesImpl {
  static_assert(EL::isList<Ind>, "Internal error of Apply!");
};

template<template<auto> class F, template<class...> class Tuple, auto t, int... es>
struct ApplyByIndicesImpl<F, EL::List<int, es...>, Tuple, t> {
  static constexpr auto value =
      Tuple<decltype(F<Get<es>(t)>::value)...>(F<Get<es>(t)>::value...);
  using value_type = std::remove_cv_t<decltype(value)>;
};

template<template<auto> class F, class Ind, template<class...> class Tuple, auto t>
constexpr auto ApplyByIndices = ApplyByIndicesImpl<F, Ind, Tuple, t>::value;
\end{cppcode}
Как обычно первая общая версия использована для проверки самого себя от внутренних ошибок.
Последний шаблон -- это сведение к вспомогательному.
А внутри имплементации я просто получаю доступ к элементу с помощью функции \verb"Get" и подставляю результат в \verb"F".
Так же мне надо явно вычислить тип этого элемента, что делается в \verb"decltype".
При этом в \verb"value_type" я добавляю \verb"std::remove_cv_t", потому что у \verb"value" тип будет обязательно с модификатором \verb"const", а мне он тут не нужен.

\subsubsection{Filter}

Так же мне понадобится вспомогательный метод для отбора элементов \verb"Tuple" по предикату.
Давайте напишу его имплементацию.
Публичная часть
\begin{cppcode}
template<template<auto> class F, auto t>
using Filter_t = typename detail::FilterImpl<F, std::remove_cv_t<decltype(t)>, t>::type;

template<template<auto> class F, auto t>
constexpr auto Filter = detail::FilterImpl<F, std::remove_cv_t<decltype(t)>, t>::value;
\end{cppcode}
Теперь имплементация, где я просто распаковываю \verb"Tuple".
\begin{cppcode}
template<template<auto> class F, class T, auto t>
struct FilterImpl {
  static_assert(TL::isList<T>, "The second element of Filter MUST be of Tuple-like type!");
};

template<template<auto> class F, class... Ts, template<class...> class Tuple, Tuple<Ts...> t>
struct FilterImpl<F, Tuple<Ts...>, t> {
  using Indices = EL::MakeList<int, 0, TL::Size<Tuple<Ts...>>>;
  static constexpr auto value = FilterByIndices<F, Indices, Tuple, t>;
  using value_type = std::remove_cv_t<decltype(value)>;
  using type = value_type;
};
\end{cppcode}
Внутри как и в случае \verb"Apply" я завожу индексы элементов и потом вызываю вспомогательный шаблон \verb"FilterByIndices", который распаковывает эту последовательность.
Вспомогательный шаблон выглядит так:
\begin{cppcode}
template<template<auto> class F, class Ind, template<class...> class Tuple, auto t>
struct FilterByIndicesImpl {
  static_assert(EL::isList<Ind>, "Internal error of Filter!");
};

template<template<auto> class F, template<class...> class Tuple, auto t, int... es>
struct FilterByIndicesImpl<F, EL::List<int, es...>, Tuple, t> {
  using DIList = EL::List<int, (F<Get<es>(t)>::value ? es : -1)...>;
  using IList = EL::EraseAll<DIList, -1>;
  using FT = ApplyByIndicesImpl<Identity, IList, Tuple, t>;
  static constexpr auto value = FT::value;
  using value_type = FT::value_type;
};
\end{cppcode}
где
\begin{cppcode}
struct Identity {
  using value_type = decltype(x);
  static constexpr value_type value = x;
};
\end{cppcode}
Давайте прокомментирую имплементацию \verb"FilterByIndicesImpl".
\begin{enumerate}
\item В первой строке я вычисляю значение предиката для всех индексов и запоминаю индекс если предикат \verb"true" и флаг \verb"-1", если \verb"false".
\begin{cppcode}[\nonumbers]
using DIList = EL::List<int, (F<Get<es>(t)>::value ? es : -1)...>;
\end{cppcode}

\item После чего я удаляю все \verb"-1"
\begin{cppcode}[\nonumbers]
using IList = EL::EraseAll<DIList, -1>;
\end{cppcode}
что оставляет только те индексы, для которых предикат верен.

\item Теперь надо отобрать элементы по индексам.
Но это можно сделать так: применим \verb"ApplyByIndicesImpl" используя тождественное отображение.
Потому что эта функция отбирала все элементы по указанным индексам и применяла функцию к каждому.
\begin{cppcode}[\nonumbers]
using FT = ApplyByIndicesImpl<Identity, IList, Tuple, t>;
\end{cppcode}

\item А дальше идет стандартное определение \verb"value" и \verb"value_type".
\end{enumerate}
